\section{Abstract}
\label{sec:abstract}

Recent literature extensively applies Mixed-Integer Linear Programming (MILP) to solve complex optimization problems in emergency management
. Within this domain, rapid and effective emergency response to urban building fires remains a critical challenge
. In these high-pressure situations, where incident commanders face dynamic and uncertain environments, Operations Research (OR) methodologies provide the necessary framework for fast, robust, and mathematically optimal decision-making
. Motivated by these operational needs, this paper introduces MORAv2, a novel adaptation of the Multicriteria Optimization-Based Resource Allocation (MORA) framework. Originally developed for large-scale wildfire suppression
, MORAv2 is re-engineered specifically for building fire emergencies.
The proposed MORAv2 model formulates indoor emergency response as an Integer Linear Programming (ILP) constrained optimization problem. Unlike traditional models that focus on macro-level station location or fleet routing
, this study focuses on micro-level tactical response actions during an active incident to maximize rescued individuals and minimize property loss
. While retaining the core lexicographic concepts and operational constraints of the original MORA framework, this research introduces a dynamic "smoke factor" to realistically capture the hazardous conditions of indoor rescue planning
. The resulting building rescue plan explicitly details which critical locations to control, optimal deployment timing, and safe routing paths, deliberately preventing responder engagement in compromised or unsurvivable situations
. Ultimately, this research bridges the critical gap between passive smart-building sensor technologies and active operational research
, providing incident commanders with a mathematically validated decision support system for both enhanced preparedness and real-time operational response.